% !TeX root = ../manual.tex
% ===========================================================================
%  Acceso al sistema
% ===========================================================================

\section{Acceso al sistema}
\label{sec:acceso}

Toda persona que use ReservaSport —usuario final o administrador— entra por
la misma pantalla de inicio de sesión. Lo que ve después de entrar depende
de su rol.

\subsection{Crear una cuenta}
\label{sec:acceso:crear-cuenta}

Si todavía no tienes una cuenta, puedes crear una desde la propia pantalla
de acceso:

\begin{enumerate}
  \item Abre la aplicación. Se muestra la pantalla \boton{Iniciar sesión} /
        \boton{Crear cuenta}.
  \item Pulsa la pestaña \boton{Crear cuenta}.
  \item Completa los tres campos:
        \begin{itemize}
          \item \campo{Usuario}: el nombre con el que iniciarás sesión.
                No puede repetirse con el de otra cuenta.
          \item \campo{Nombre completo}: cómo quieres que te llamen dentro
                de la aplicación.
          \item \campo{Contraseña}: mínimo 6 caracteres.
        \end{itemize}
  \item Pulsa \boton{Crear cuenta}.
\end{enumerate}

\captura{registro-form-vacio}{1}{Formulario de creación de cuenta.}{registro-vacio}

\captura{registro-form-completo}{1}{Formulario completo, listo para enviar.}{registro-completo}

\nota{Las cuentas que se crean desde aquí siempre quedan como
\textbf{usuario final}. Si necesitas una cuenta de administrador, debe
crearla o ascenderla el club desde Gestión de usuarios (véase
\cref{sec:admin:usuarios}).}

Al crear la cuenta, el sistema inicia sesión automáticamente y te lleva a
\boton{Mis Reservas}. Como es una cuenta nueva, todavía no tiene ninguna
reserva.

\captura{registro-exito}{1}{Cuenta creada: sesión iniciada automáticamente, sin reservas todavía.}{registro-exito}

\subsection{Iniciar sesión}
\label{sec:acceso:login}

\begin{enumerate}
  \item En la pestaña \boton{Iniciar sesión}, escribe tu \campo{Usuario} y
        tu \campo{Contraseña}.
  \item Pulsa \boton{Iniciar Sesión}.
\end{enumerate}

\captura{usuario-login}{1}{Inicio de sesión.}{login-usuario}

Si tus datos son correctos, entras directamente a la sección que
corresponde a tu rol: \boton{Mis Reservas} si eres usuario final, o el
\boton{Dashboard} de administración si eres administrador.

\subsubsection{Mensajes al iniciar sesión}

El sistema distingue dos motivos de rechazo y te lo dice explícitamente,
para que sepas qué hacer:

\captura{login-error-credenciales}{1}{Usuario o contraseña incorrectos.}{login-error-credenciales}

\aviso{Si ves \textit{«Usuario o contraseña incorrectos»}, revisa que no
haya espacios de más y que el bloqueo de mayúsculas esté desactivado. Por
seguridad, el sistema no indica cuál de los dos datos falló.}

\captura{login-error-inactivo}{1}{Cuenta inactiva: se indica el motivo y a quién contactar.}{login-error-inactivo}

\aviso{Si ves \textit{«La cuenta está inactiva»}, tu cuenta existe pero un
administrador la desactivó. Contacta al club para que la reactive desde
Gestión de usuarios (\cref{sec:admin:usuarios}); tus reservas anteriores no
se pierden mientras tanto.}

\subsection{Mi Perfil}
\label{sec:acceso:perfil}

Desde el menú lateral, la opción \boton{Perfil} muestra los datos de tu
cuenta: nombre, nombre de usuario, rol e identificador interno. Es la
misma pantalla para usuario final y administrador; solo cambia el rol que
se muestra. El menú lateral permanece visible mientras consultas tu
perfil, igual que en el resto de la aplicación.

\captura{usuario-perfil}{1}{Perfil de un usuario final. El menú lateral se mantiene visible.}{usuario-perfil}

\captura{admin-perfil}{1}{Perfil de un administrador. El menú lateral se mantiene visible.}{admin-perfil}

\subsection{Cerrar sesión}

Pulsa \boton{Cerrar Sesión}, disponible tanto en el menú lateral como en la
propia pantalla de \boton{Perfil}. La sesión se cierra de inmediato y
vuelves a la pantalla de acceso.

\subsubsection{Cierre de sesión por inactividad}

Si pasas un minuto sin mover el mouse, escribir ni desplazarte, la
aplicación muestra un aviso con una cuenta regresiva de 30 segundos antes
de cerrar la sesión automáticamente.

\captura{sesion-inactividad}{1}{Aviso de inactividad, con la cuenta regresiva antes del cierre automático.}{sesion-inactividad}

\begin{enumerate}
  \item Pulsa \boton{Estoy Aquí} para seguir trabajando; la cuenta
        regresiva se cancela y el conteo de inactividad vuelve a empezar.
  \item Si no respondes, la sesión se cierra sola y vuelves a la pantalla
        de acceso, igual que si hubieras pulsado \boton{Cerrar Sesión}.
\end{enumerate}

\nota{Es una medida de seguridad: evita que una sesión quede abierta en un
equipo compartido. No se pierde ningún dato por el cierre: cualquier
reserva o cambio ya guardado permanece como estaba.}
